\begin{PosterAppCard}{WayBuddy}
\PosterStatus{Experimenteller Demonstrator · offline}\hfill\PosterAppLogo{assets/figures/way-buddy-logo.webp}

\vspace{4mm}
\PosterLead{Forschungsdemo für lokale Freiflächen-Erkennung.}

\PosterAppImage{assets/figures/way-buddy-screenshot.webp}{WayBuddy Screenshot: begehbare Fläche und Richtungsfeedback}

\PosterAppBullets
  {Kamerabild lokal auf dem Smartphone analysieren}
  {Begehbare Flächen per DNN segmentieren}
  {Richtung und Unsicherheit akustisch signalisieren}
  {Machbarkeit und Grenzen transparent demonstrieren}

\vspace{3mm}
\PosterTag{Offline-DNN}\hspace{1mm}\PosterTag{Freifläche}\hspace{1mm}\PosterTag{Audio}\hspace{1mm}\PosterTag{UNKNOWN}

\vspace{5mm}
\PosterSmallNote{Wichtig: WayBuddy ist kein fertiges Blinden-Navigationssystem und ersetzt keine etablierten Mobilitätshilfen. Geeignet: \texttt{screenshot-01} oder \texttt{screenshot-02}.}
\end{PosterAppCard}
